\documentclass[10pt,a4paper]{amsart}
\usepackage[margin=19mm]{geometry}
\usepackage{amsmath,amssymb,mathtools}
\usepackage[scr=rsfs,cal=euler]{mathalfa}
\usepackage{xfrac}
\usepackage[unicode,hidelinks,bookmarks=false]{hyperref}
\newcommand{\ie}{i.e.\ }
\newcommand{\cf}{cf.\ }
\newcommand{\resp}{respectively\ }
\newcommand{\Subsection}{\subsection}
\providecommand{\nobreakdash}{\nobreak\hbox{-}}
\setlength{\emergencystretch}{2em}
\multlinegap0pt
\usepackage{caption}
\usepackage[shortlabels]{enumitem}
\usepackage{bbm}
\usepackage{amssymb}
\usepackage{mathtools}
\usepackage{booktabs}
\usepackage{algorithm}\usepackage{algpseudocode}
\newtheorem{assumption}{Assumption}
\newtheorem{lemma}{Lemma}
\newcommand{\Cov}{{\rm Cov}}
\newcommand{\E}{\mathbb{E}}
\newcommand{\R}{\mathbb{R}}
\newcommand{\Range}{{\rm Range}}
\newcommand{\Sym}{{\rm \mathrm{Sym}}}
\newcommand{\norm}[1]{\left\|{#1}\right\|} % A norm with 1 argument
\title{A Theory of Feature Learning in Kernel Models}
\author{Yunlu Chen, Yang Li, Keli Liu,, Feng Ruan}
\date{Source: arXiv:2310.11736v4 (CC BY 4.0)}
\begin{document}
\maketitle

\noindent\emph{Source and licence.} A Theory of Feature Learning in Kernel Models. Version arXiv:2310.11736v4 (CC BY 4.0), distributed by arXiv under the Creative Commons Attribution 4.0 International licence, http://creativecommons.org/licenses/by/4.0/, as recorded in the arXiv licence metadata for this exact version and shown on the abstract page https://arxiv.org/abs/2310.11736v4. Full licence notice: \texttt{licenses/arxiv-2310.11736-CC-BY-4.0.txt}.\par\smallskip

\noindent\textit{Editorial scope.} Counted unit: the article's lemma lemma:projection-reduce-function-value (source0.tex lines 1843--1850). The article's complete original proof of it is delivered with it (source0.tex lines 1852--1919, one \texttt{proof} environment of 68 lines). No mathematics is written, repaired or re-derived: the statement and the proof are the article's own lines, in the article's own notation, and no retained line is substituted or re-flowed.  Retained uncounted as the source-local context the proof needs: lines 607--611; lines 949--952.  Nothing was omitted from any retained span, so the package carries no omission record.  No retained line contains a citation, so the wrapper carries no bibliography and no bibliography entry.  No retained line contains a figure environment, an image-inclusion command or drawing code, so no image is delivered and the package declares no asset dependency.  Editorial formatting only, in addition: an AMS \texttt{amsart} preamble that restates the article's own declarations and macros used by the retained lines, namely the environments \texttt{assumption}, \texttt{lemma} on separate counters, together with the standard packages those lines need.  The retained environments print in this document as Assumption 1, Lemma 1 in order of appearance, whereas the article prints the same items with the numbering of its own full document.  The complete native source of the article as distributed in the arXiv e-print for this exact version is bundled unchanged as \texttt{sources/148657/source0.tex}.
\begin{assumption}
\label{assumption:independence}
The random vectors $\Pi_S X$ and $\Pi_{S^\perp} X$ are probabilistically independent.  
Moreover, the random vector $\Pi_{S^\perp} X$ is not degenerate (i.e., $\Cov(\Pi_{S^\perp} X)$ is not zero). 
\end{assumption} 

\begin{lemma}
\label{lemma:continuous-embedding}
For every $f \in H$, we have $\norm{f}_{L_\infty} \le \norm{f}_H$, and $f$ has a continuous representative, and $\lim_{x \to \infty} f(x) = 0$.
\end{lemma} 

\begin{lemma}
\label{lemma:projection-reduce-function-value}
For every $\lambda > 0$
\begin{equation*}
	J(\Sigma, \lambda) \ge J(\Pi_S \Sigma \Pi_S, \lambda),~~~\forall \Sigma \in \Sym_+^d.
\end{equation*} 
The inequality is strict when $\Range(\Sigma) \not \subset S$ and $J(\Sigma, \lambda) < \E[Y^2]$.
\end{lemma} 

\begin{proof} 
Fix $U$ such that $\Sigma = U^\top U$. 
For every $\xi$ in $\R^d$, we define the translated function $\tau_\xi f$ by
\begin{equation*}
	(\tau_\xi f)(x):= f(x+\xi),~~~\forall x \in \R^d.
\end{equation*}
By the translation-invariance of the RKHS $H$, we have $\norm{\tau_\xi f}_H = \norm{f}_H$. 
By the minimizing property of $\mathcal{J}$, we deduce for every vector $\xi$ in $\R^d$, and $f \in H$ that 
\begin{equation*}
\begin{aligned}
  \E[(Y - f(U\Pi_S X + \xi))^2] + \lambda \|f\|_H^2
   &= \E[(Y - (\tau_\xi f)(U\Pi_S X))^2] + \lambda \|\tau_\xi f\|_H^2 \\
   &= \mathcal{I}(\tau_\xi f, U\Pi_S, \lambda)
   \;\ge\; \mathcal{J}(U\Pi_S, \lambda).
\end{aligned}
\end{equation*}
Define $c := \E[(Y - \E[Y|X])^2]$. Since $\E[Y|X] = \E[Y|\Pi_S X]$, we can decompose the squared error as
\[
\E[(Y - f(U \Pi_{S} X + \xi))^2] = \E\left[(\E[Y|\Pi_{S} X] - f(U \Pi_{S} X + \xi))^2\right] + c.
\]
Substituting into the previous inequality gives
\[
\E\left[(\E[Y|\Pi_{S} X] - f(U \Pi_{S} X + \xi))^2\right] + \lambda \|f\|_H^2 + c \ge \mathcal{J}(U\Pi_S, \lambda).
\]

By Assumption~\ref{assumption:independence}, $\Pi_{S} X$ and $\Pi_{S^\perp} X$ are independent. Setting $\xi := U \Pi_{S^\perp} X$, we get
\[
  \E\left[(\E[Y|\Pi_{S} X] - f(U \Pi_{S} X + U\Pi_{S^\perp} X))^2 \mid  \Pi_{S^\perp} X\right] + \lambda \|f\|_H^2 + c \ge \mathcal{J}(U\Pi_S, \lambda).
\]
Taking expectation on both sides, and using $I = \Pi_{S} + \Pi_{S^\perp}$, we get 
\begin{equation*}
\begin{split} 
	\mathcal{I}(f, U, \lambda) 
		&=  \E[(Y - f( U X))^2] + \lambda \norm{f}_H^2  \\
		&= \E[(Y - f( U \Pi_{S} X + U \Pi_{S^\perp} X))^2] + \lambda \norm{f}_H^2 \\
		&=   \E\left[(\E[Y|\Pi_{S} X] - f(U \Pi_{S} X + U \Pi_{S^\perp} X))^2\right] + \lambda \|f\|_H^2 + c 
			\ge  \mathcal{J}(U\Pi_S, \lambda).
\end{split} 
\end{equation*}
By minimizing over $f \in H$, we get 
$\mathcal{J}(U, \lambda)= \min_f \mathcal{I}(f, U, \lambda)  \ge \mathcal{J}(U\Pi_S, \lambda)$. 
Since $\Sigma = U^\top U$, 
\begin{equation*} J(\Sigma,\lambda) = \mathcal{J}(U,\lambda) \ge \mathcal{J}(\Pi_S U,\lambda)
  = J(\Pi_S \Sigma \Pi_S,\lambda).
 \end{equation*} 
%$(U\Pi_S)^\top (U\Pi_S) = \Pi_S \Sigma \Pi_S$, so by the
%definition of $J$,
%\[
%  J(\Sigma,\lambda) = \mathcal{J}(U,\lambda) \ge \mathcal{J}(\Pi_S U,\lambda)
%  = J(\Pi_S \Sigma \Pi_S,\lambda).
%\]
%This proves the claimed inequality.

%Recall that for each $\beta$, the functional $f \mapsto \mathcal{I}(f, \beta, \lambda)$ admits a unique minimizer, denoted $f_{\beta, \lambda}$. 
Our derivation shows equality 
$J(\Sigma, \lambda) = J(\Pi_{S} \Sigma \Pi_S, \lambda)$ holds
only if the minimizer $f_{U, \lambda}$ obeys 
\begin{equation*}
(\tau_{\xi}) f_{U, \lambda} = (\tau_{\xi'}) f_{U, \lambda}~~\text{for almost all}~\xi= U\Pi_{S^\perp} X, \, \xi' = U\Pi_{S^\perp}X'
\end{equation*}
where $X, X'$ are independent copies. If $\Pi_{S^\perp}\Sigma\Pi_{S^\perp}\neq 0$, then $U\Pi_{S^\perp}\neq 0$, and thus $\operatorname{Cov}(U\Pi_{S^\perp}X)\neq 0$. Therefore, with positive probability $\xi\neq \xi'$. Fix such a realization and set $z=\xi-\xi'\neq 0$. Then
\begin{equation*}
	f_{U, \lambda}(x + z)  = f_{U, \lambda}(x)~~~\forall x\in \R^d.
\end{equation*}
However, since $f_{U, \lambda} \in H$, it satisfies $\lim_{x \to \infty} f_{U, \lambda}(x) = 0$ 
by Lemma~\ref{lemma:continuous-embedding}. In turn, this implies that $f_{U, \lambda}$ must be identically zero. 
Consequently, $\mathcal{J}(U, \lambda) = \E[Y^2]$, and therefore $J(\Sigma, \lambda) = \E[Y^2]$.
\end{proof} 

\end{document}
